\documentclass[tikz,border=10pt]{standalone}
\usetikzlibrary{positioning, arrows.meta, calc, babel, fit, backgrounds}
\usepackage[utf8]{inputenc}
\usepackage[T1]{fontenc}
\usepackage{lmodern}
\usepackage[spanish,es-tabla]{babel}
\usepackage{amsmath, amssymb, bm}

\begin{document}
\begin{tikzpicture}[
  align=center,
  font=\sffamily\small,
  % Estilos de bloques
  phasebox/.style={draw, rectangle, rounded corners=3pt, minimum width=2.5cm, minimum height=1.1cm, thick, fill=white},
  artbox/.style={draw, rectangle, dashed, rounded corners=2pt, minimum width=2.2cm, minimum height=0.9cm, fill=yellow!8!orange!2!white, draw=orange!60!black, font=\sffamily\scriptsize, thick},
  groupbox/.style={draw, rectangle, rounded corners=5pt, line width=1.0pt},
  % Estilos de flechas
  line/.style={draw, -{Stealth[scale=1.2]}, thick},
  dashline/.style={draw, -{Stealth[scale=1.2]}, thick, dashed},
  dashseg/.style={draw, thick, dashed}
]

  % ==========================================
  % FILA 1: ETAPA I (y = 0.0)
  % ==========================================
  \node[phasebox] (f1) at (0,0) {\textbf{Fase 1}\\Sanidad y Tests\\(\texttt{pytest})};
  \node[phasebox] (f2) at (3.6,0) {\textbf{Fase 2}\\Dataset Clásico\\Inicial};
  \node[phasebox] (f3) at (7.2,0) {\textbf{Fase 3}\\Diagnóstico y\\Ajuste PID};
  \node[phasebox] (f4) at (10.8,0) {\textbf{Fase 4}\\Baseline Clásico\\con PIDs Congelados};

  \draw[line] (f1) -- (f2);
  \draw[line] (f2) -- (f3);
  \draw[line] (f3) -- (f4);

  % Artefacto de PIDs congelados - centrado en X=9.0, y=-1.55
  \node[artbox] (art_pids) at (9.0,-1.55) {YAMLs de PIDs\\congelados por\\escenario};
  
  \draw[dashline] (f3.south) |- (art_pids.west);

  % ==========================================
  % FILA 2: ETAPA II (Fase 5 en y=-3.9, Fase 6 en y=-5.5, Fase 7 y art_weights en y=-4.7)
  % ==========================================
  \node[phasebox] (f5) at (2.0,-4.1) {\textbf{Fase 5}\\Dataset Fuerza\\Externa (Imitación)};
  \node[phasebox] (f6) at (2.0,-5.7) {\textbf{Fase 6}\\Dataset Ganancias\\de Posición};
  \node[phasebox] (f7) at (7.2,-4.7) {\textbf{Fase 7}\\Entrenamiento\\MLP / GRU / LSTM};
  \node[artbox] (art_weights) at (10.8,-4.7) {Checkpoints de\\modelos entrenados\\(\texttt{best.pt})};

  % Entradas de la Etapa I a las Fases 5 y 6 (cableado paralelo continuo/discontinuo)
  \coordinate (in_cont) at (-0.5,-3.25);
  \coordinate (in_dash) at (0,-3.1);
  
  % Conexión del YAML (art_pids) a Fase 4 (entrada en X=10.5)
  \draw[dashline] (art_pids.east) -- (10.5,-1.55) -- (10.5,-0.55);

  % Canal de datos continuo: sale de f4.south desplazado a la derecha (en X=11.1), baja y va a la izquierda
  \draw[thick] ($(f4.south)+(0.3,0)$) -- (11.1,-3.25) -- (in_cont);
  \draw[line] (in_cont) |- ($(f5.west)+(0,0.15)$);
  \draw[line] (in_cont) |- ($(f6.west)+(0,0.15)$);

  % Canal de parámetros discontinuo: derivación en X=10.5, baja y va a la izquierda
  \draw[thick, dashed] (10.5,-1.55) -- (10.5,-3.1) -- (in_dash);
  \draw[dashline] (in_dash) |- ($(f5.west)+(0,-0.15)$);
  \draw[dashline] (in_dash) |- ($(f6.west)+(0,-0.15)$);

  % Conexión limpia en "Y" de F5/F6 hacia F7 (flechas más directas al estar desplazadas a la izquierda)
  \draw[line] (f5.east) -- ++(0.8,0) |- ($(f7.west)+(0,0.25)$);
  \draw[line] (f6.east) -- ++(0.8,0) |- ($(f7.west)+(0,-0.25)$);
  
  \draw[line] (f7.east) -- (art_weights.west);

  % ==========================================
  % FILA 3: ETAPA III (y = -9.0, Panel más alto)
  % ==========================================
  \node[phasebox] (f8) at (0,-9.0) {\textbf{Fase 8}\\Evaluación\\Supervisada};
  \node[phasebox] (f9) at (3.6,-9.0) {\textbf{Fase 9}\\Simulación bucle\\cerrado (Test ID)};
  \node[phasebox] (f10) at (7.2,-9.0) {\textbf{Fase 10}\\Simulación bucle\\cerrado (OOD)};
  \node[phasebox] (f11) at (10.8,-9.0) {\textbf{Fase 11}\\Consolidación y\\tablas LaTeX};
  \node[artbox] (art_latex) at (10.8,-10.65) {Tablas comparativas\\de rendimiento en\\LaTeX (\texttt{results/})};

  % Bus de checkpoints: baja de art_weights y corre horizontal a Y=-7.8 (por debajo del título)
  \draw[thick, dashed] (art_weights.south) -- (10.8,-7.8) -- (0,-7.8);
  
  \draw[dashline] (0,-7.8) -- (f8.north);
  \draw[dashline] (3.6,-7.8) -- (f9.north);
  \draw[dashline] (7.2,-7.8) -- (f10.north);

  \draw[line] (f8) -- (f9);
  \draw[line] (f9) -- (f10);
  \draw[line] (f10) -- (f11);
  
  \draw[dashline] (f11.south) -- (art_latex.north);

  % ==========================================
  % LÍNEA BASE CLÁSICA (Extremo derecho, exterior)
  % ==========================================
  \draw[line, draw=black!30] (f4.east) -- (12.5,0) -- (12.5,-9.0) -- (f11.east);
  \node[right, font=\sffamily\scriptsize\itshape, text=black!50] at (12.6,-4.5) {Línea base clásica};

  % ==========================================
  % CAPA DE FONDO: PANELES DE ETAPA (y sep=0.2cm)
  % ==========================================
  \begin{scope}[on background layer]
    % Panel Etapa I
    \draw[groupbox, fill=blue!5!gray!2, draw=blue!40!gray!50] (-1.4, 1.4) rectangle (12.2, -2.2);
    % Panel Etapa II
    \draw[groupbox, fill=red!5!gray!2, draw=red!40!gray!50] (-1.4, -2.4) rectangle (12.2, -6.6);
    % Panel Etapa III (Más alto)
    \draw[groupbox, fill=green!5!gray!2, draw=green!40!gray!50] (-1.4, -6.8) rectangle (12.2, -11.3);
  \end{scope}

  % Títulos de las Etapas colocados en la esquina superior izquierda
  \node[anchor=north west, font=\sffamily\bfseries\footnotesize, color=blue!60!black] at (-1.25, 1.25) {ETAPA I: SINTONIZACIÓN Y LÍNEA DE BASE CLÁSICA};
  \node[anchor=north west, font=\sffamily\bfseries\footnotesize, color=red!60!black] at (-1.25, -2.55) {ETAPA II: GENERACIÓN DE DEMOSTRACIONES Y ENTRENAMIENTO NEURONAL};
  \node[anchor=north west, font=\sffamily\bfseries\footnotesize, color=green!50!black] at (-1.25, -6.95) {ETAPA III: EVALUACIÓN Y CONSOLIDACIÓN DE RESULTADOS};

\end{tikzpicture}
\end{document}